\documentclass[11pt]{article}
\usepackage[margin=25mm]{geometry}
\usepackage{fontspec}
\setmainfont{Latin Modern Roman}
\usepackage{amsmath,amssymb,mathtools}
\usepackage{unicode-math}
\setmathfont{Latin Modern Math}
\usepackage{xurl,hyperref}
\hypersetup{colorlinks=true,urlcolor=blue}
\setcounter{secnumdepth}{0}
\setlength{\parindent}{0pt}
\setlength{\parskip}{5pt}
\setlength{\emergencystretch}{3em}
\newcommand{\sref}[2]{\hyperlink{source:#1:#2}{[S#2]}}
\newcommand{\cataloguescope}[1]{\paragraph{Recorded target.}#1}
\begin{document}
\section{Nonseparating planar continuum}\label{statement}

\subsection{Definitions and mathematical statement}

\paragraph{Shared notation.}$\mathbb N=\{1,2,\ldots\}$, and $\mathbb Z,\mathbb Q,\mathbb R,\mathbb C$ denote the integers, rationals, reals and complex numbers. Any other conventions are specified locally.


A nonempty compact set $K\subseteq\mathbb R^2$ is path connected if any two points are joined by a continuous map $[0,1]\to K$. It is nonseparating if $\mathbb R^2\setminus K$ is connected. It has the fixed-point property if every continuous map $f:K\to K$ fixes some point. Continuity uses the subspace topology.
\paragraph{Statement recorded in the supplied source.}
Does every nonempty path-connected nonseparating compact set $K\subseteq\mathbb R^2$ have the fixed-point property? Explicitly, does
\[
 \forall f\in C(K,K)\quad\exists x\in K:\quad f(x)=x
\]
follow from those hypotheses? \sref{statement}{2}
\subsection{Short English statement}
Must a continuous self-map of a path-connected planar compact set with connected complement have a fixed point?
\subsection{Sources}
\begin{description}
\item[\hypertarget{source:statement:1}{S1}] ULAM.AI and the UnsolvedMath Contributors, UnsolvedMath: A Curated Collection of Open Mathematics Problems, catalogue entry. CC BY 4.0. \url{https://huggingface.co/datasets/ulamai/UnsolvedMath}.
\item[\hypertarget{source:statement:2}{S2}] Open Problem Garden, Nonseparating planar continuum, catalogue entry. Original problem statement, full mathematical discussion, bibliography and comments. \url{https://www.openproblemgarden.org/op/nonseparating_planar_continuum}.
\end{description}
\label{end:statement}
\end{document}
